% Chapter VII, second half: §72–74 + Bibliographical Notes
% The Stability of Couette Flow (continued)
% Continues from chapter_7_part1.tex (which ends at eq:7-275)

%% §72 %%
\section{On the principle of the exchange of stabilities}\label{sec:7-72}

As we have remarked in \S\,71 and as we shall see in detail in \S\,74,
experiments confirm that the instability of the Couette flow sets in as a
secondary stationary flow.  Without attempting a full theoretical
justification of the principle of the exchange of stabilities, we shall give
some reasons why we may expect its validity in this instance.  For this
purpose, we return to equations~\eqref{eq:7-196} and~\eqref{eq:7-197} in which the time
constant~$\sigma$ is still retained.  Translating the origin of the coordinate
system to be midway between the two cylinders and replacing~$\mathfrak{u}$ by
$\mathfrak{u}(1+\mu)/2$, we obtain (cf.\ equations~\eqref{eq:7-201}--\eqref{eq:7-204})
\begin{equation}
\label{eq:7-276}
(D^2-a^2)(D^2_*-a^2-\sigma)\mathfrak{u} = (1+\varepsilon x)v
\end{equation}
and
\begin{equation}
\label{eq:7-277}
(D^2-a^2-\sigma)v = -\mathscr{T}a^2\mathfrak{u},
\end{equation}
where
\begin{equation}
\label{eq:7-278}
\mathscr{T} = \tfrac{1}{4}(1+\mu)T \quad \text{and} \quad
\varepsilon = -a\,\frac{1-\mu}{1+\mu}.
\end{equation}
And the boundary conditions are the same as hitherto.

Consider first the case $\mu > 0$.  As we have seen in \S\,71\,(d), in this case,
when $\sigma = 0$, the lowest characteristic values of the problem hardly
depend on the term linear in~$x$ on the right-hand side of equation~\eqref{eq:7-276}.
In fact, to order~$\varepsilon$, the characteristic values are the same as when $\varepsilon = 0$;
and, moreover, on account of the wide spacing of the characteristic
values of the `unperturbed problem', the coefficient of~$\varepsilon^2$ in a perturbation
series for~$\mathscr{T}$ is very small (see equation~\eqref{eq:7-263}).  It would, therefore,
appear that even when $\sigma \ne 0$, we can to the \emph{first} order in~$\varepsilon$ ignore the term
in~$\varepsilon x$ on the right-hand side of equation~\eqref{eq:7-276}.  We shall, then, be left with
\begin{equation}
\label{eq:7-279}
(D^2-a^2)(D^2-a^2-\sigma)\mathfrak{u} = v,
\end{equation}
and
\begin{equation}
\label{eq:7-280}
(D^2-a^2-\sigma)v = -\mathscr{T}a^2\mathfrak{u},
\end{equation}
together with the boundary conditions
\begin{equation}
\label{eq:7-281}
v = \mathfrak{u} = D\mathfrak{u} = 0 \quad \text{for } x = \pm\tfrac{1}{2}.
\end{equation}

Now multiply equation~\eqref{eq:7-279} by~$\mathfrak{u}^*$ and integrate over the range of~$x$.
We obtain by an integration by parts
\begin{equation}
\label{eq:7-282}
\int_{-\frac{1}{2}}^{+\frac{1}{2}} v\mathfrak{u}^*\,dx
= \int_{-\frac{1}{2}}^{+\frac{1}{2}}
  \bigl\{\mathfrak{u}^*\!\bigl[(D^2-a^2)^2 - \sigma(D^2-a^2)\bigr]\mathfrak{u}\bigr\}\,dx
= \int_{-\frac{1}{2}}^{+\frac{1}{2}}
  \bigl\{|(D\mathfrak{u}|^2+a^2|\mathfrak{u}|^2\bigr)^2\bigr\}\,dx
+ \sigma\!\int_{-\frac{1}{2}}^{+\frac{1}{2}}
  \bigl(|D\mathfrak{u}|^2+a^2|\mathfrak{u}|^2\bigr)\,dx.
\end{equation}
On the other hand, from equation~\eqref{eq:7-280} we have
\begin{equation}
\label{eq:7-283}
-\mathscr{T}a^2\!\int_{-\frac{1}{2}}^{+\frac{1}{2}} v\mathfrak{u}^*\,dx
= \int_{-\frac{1}{2}}^{+\frac{1}{2}}
  v(D^2-a^2-\sigma)v^*\,v^*\,dx
= -\int_{-\frac{1}{2}}^{+\frac{1}{2}}
  \bigl\{|Dv|^2+(a^2+\sigma^*)|\,v|^2\bigr\}\,dx.
\end{equation}
Combining equations~\eqref{eq:7-282} and~\eqref{eq:7-283}, we obtain
\begin{multline}
\label{eq:7-284}
\mathscr{T}a^2\!\int_{-\frac{1}{2}}^{+\frac{1}{2}}
  \bigl\{|(D^2-a^2)\mathfrak{u}|^2\,dx
  + \sigma\!\int_{-\frac{1}{2}}^{+\frac{1}{2}}
    (|D\mathfrak{u}|^2+a^2|\mathfrak{u}|^2)\,dx\bigr\} \\
= \int_{-\frac{1}{2}}^{+\frac{1}{2}}
  \bigl(|Dv|^2+a^2|v|^2\bigr)\,dx
  + \sigma^*\!\int_{-\frac{1}{2}}^{+\frac{1}{2}} |v|^2\,dx.
\end{multline}
On equating the imaginary parts of equation~\eqref{eq:7-284}, we obtain
\begin{equation}
\label{eq:7-285}
\operatorname{im}(\sigma)\!\Biggl\{
\mathscr{T}a^2\!\int_{-\frac{1}{2}}^{+\frac{1}{2}}
  \bigl(|D\mathfrak{u}|^2+a^2|\mathfrak{u}|^2\bigr)\,dx
  + \int_{-\frac{1}{2}}^{+\frac{1}{2}} |v|^2\,dx\Biggr\}
= 0.
\end{equation}
From this equation it follows that
\begin{equation}
\label{eq:7-286}
\operatorname{im}(\sigma) = 0 \quad \text{when } \mathscr{T} > 0;
\end{equation}
and we conclude that the principle of the exchange of stabilities is valid.
In view of the fact that for $\mu > 0$ the first-order formula~\eqref{eq:7-231} gives
values which differ from the results of more exact calculations by less
than one per cent, it would appear that the foregoing arguments effectively exclude the occurrence of overstability for $\mu > 0$.

As~$\mu$ turns negative, the replacement of equations~\eqref{eq:7-276} and~\eqref{eq:7-277} by
\eqref{eq:7-279} and~\eqref{eq:7-280} will lead to rapidly increasing errors, and the possibility
of overstability occurring cannot be excluded in the same way.  On the
other hand, from the discussion in \S\,71\,(e), it would appear that for
$(1-\mu) \to \infty$ the basic physical phenomenon is not very different from
what happens when $\mu = 0$; the only difference is that it all happens in
the layers immediately adjoining the inner cylinder; so that in this case
also, overstability would seem unlikely.

While the foregoing arguments are not conclusive, it is not, in principle,
difficult to settle the question by actual computation.  Thus, the characteristic value problem presented by equations~\eqref{eq:7-196},~\eqref{eq:7-197}, and~\eqref{eq:7-200}\footnote{We have replaced~$\sigma$ in equations~\eqref{eq:7-196} and~\eqref{eq:7-197} by~$i\sigma$ to emphasize that we are now primarily interested in marginal states which are strictly oscillatory.}
can be solved by the method described in \S\,71\,(a).  We find that equation
\eqref{eq:7-213} is replaced by
\begin{multline}
\label{eq:7-287}
\Biggl[\frac{a(m^2\pi^2+a^2+b^2)(2m^2\pi^2+a^2+b^2)}{(n^2\pi^2+a^2)(n^2\pi^2+a^2+b^2)(m^2\pi^2+a^2+b^2)}(-1)^{m+n}-1\Biggr] - {} \\
\frac{mn\pi^2(a^2-b^2)}{(n^2+a^2)(m^2\pi^2+a^2+b^2)}\bigl\{(\sigma\sinh a - a\sinh b)[(-1)^{m+1}+(1+\alpha)(-1)^{m+1}] + {} \\
{}+(\alpha\cosh b\sinh a - a\sinh b\cosh a)[1+(1+\alpha)(-1)^{m+n}] + {} \\
\frac{2\alpha(2m^2\pi^2+a^2+b^2)}{(m^2\pi^2+a^2)(m^2\pi^2+b^2)}
  \bigl[ab(\cosh a - \cosh b)[(-1)^{m+1}-(-1)^{m+1}] + {} \\
{}+\tfrac{1}{2}(a^2-b^2)\sinh a\sinh b[(-1)^{m+n}-1]\bigr]\Biggr\} + {} \\
{}+\tfrac{1}{2}\delta_{m,n} + \alpha X_{m,n}
  - \tfrac{1}{4}(n^2\pi^2+a^2)(n^2\pi^2+a^2+b^2)\frac{2m\pi}{|_{v_n}|}
= 0,
\end{multline}
where
\begin{equation}
\label{eq:7-288}
b^2 = a^2 + i\sigma,
\end{equation}
\begin{equation}
\label{eq:7-289}
\Delta = 2ab(1-\cosh a\cosh b) + (a^2+b^2)\sinh a\sinh b,
\end{equation}
and
\begin{equation}
\label{eq:7-290}
X_{m,n} = \begin{cases}
0 & \text{if } m+n \text{ is even and } m \ne n, \\
\frac{1}{2} & \text{if } m = n, \\
\dfrac{4mn}{n^2-m^2}\!\left(
  \dfrac{2m^2n^2+a^2+b^2}{(m^2\pi^2+a^2)(m^2\pi^2+b^2)}
  - \dfrac{1}{n^2(n^2-m^2)}\right)
& \text{if } m+n \text{ is odd.}
\end{cases}
\end{equation}
And the question is: Is there a real~$\sigma$ for which (for some given~$a$) the
characteristic root of the secular equation~\eqref{eq:7-287} is real?  If the answer is
in the affirmative, then overstability is possible.

A detailed examination of the roots of equation~\eqref{eq:7-287} has not been
undertaken.  It would be particularly worthwhile to explore the case
$\mu = -1$.

%% §73 %%
\section{The solution for a wide gap when the marginal state is stationary}\label{sec:7-73}

We now return to a consideration of the equations~\eqref{eq:7-167}--\eqref{eq:7-170} without
making any assumptions about the relative width of the gap $R_2-R_1$
between the cylinders.  But we shall continue to restrict ourselves to the
case when the onset of instability is as a stationary pattern of secondary
flow.  The relevant equations then are
\begin{equation}
\label{eq:7-291}
(DD_*-a^2)^2\mathfrak{u} = -Ta^2\!\left(\tfrac{1}{2}-\kappa\right)v
\end{equation}
and
\begin{equation}
\label{eq:7-292}
(DD_*-a^2)v = \mathfrak{u}.
\end{equation}
Eliminating~$\mathfrak{u}$ between equations~\eqref{eq:7-291} and~\eqref{eq:7-292}, we obtain
\begin{equation}
\label{eq:7-293}
(DD_*-a^2)^3 v = -Ta^2\!\left(\tfrac{1}{2}-\kappa\right)v;
\end{equation}
and the appropriate boundary conditions are
\begin{equation}
\label{eq:7-294}
v = (DD_*-a^2)v = D(DD_*-a^2)v = 0
\quad\text{for } r = 1 \text{ and } \eta.
\end{equation}

A solution of equation~\eqref{eq:7-293} patterned after the method described in
\S\,71\,(a) proves impracticable.  An alternative method based on expansions in terms of orthogonal functions satisfying four boundary conditions
appears well suited to the problem at hand.  The method is the following.

Letting
\begin{equation}
\label{eq:7-295}
G = (DD_*-a^2)v,
\end{equation}
we observe that the boundary conditions~\eqref{eq:7-294} require that both~$G$
and its derivative vanish at $r = 1$ and~$\eta$.  Accordingly, we may expand
$G$ in terms of the set of orthogonal functions, $\mathscr{G}_i(\alpha_i r)$, determined by the
characteristic value problem specified by the equation
\begin{equation}
\label{eq:7-296}
(DD_*)^2 y = \left(\frac{d^2}{dr^2}+\frac{1}{r}\frac{d}{dr}-\frac{1}{r^2}\right)^2 y_i = \alpha_i^4 y,
\end{equation}
and the boundary conditions
\begin{equation}
\label{eq:7-297}
y = 0 \quad\text{and}\quad dy/dr = 0 \quad\text{for } r = 1 \text{ and } \eta.
\end{equation}
The required characteristic functions are expressible as linear combinations of the Bessel functions $J_1$, $Y_1$, $I_1$, and~$K_1$ in the form
\begin{equation}
\label{eq:7-298}
\mathscr{G}_i(\alpha_i r) = A_i J_1(\alpha_i r) + B_i Y_1(\alpha_i r) + C_i I_1(\alpha_i r)
  + D_i K_1(\alpha_i r),
\end{equation}
where~$\alpha_i$ is a root of a certain transcendental equation (given in Appendix
V, equation~(31)) and $A_i$, $B_i$, $C_i$, and~$D_i$ are constants determined apart
from an arbitrary constant of proportionality.

The basic idea, then, is to expand~$G$ in terms of the functions
$\mathscr{G}_i(\alpha_i r)$.  Thus, we assume
\begin{equation}
\label{eq:7-299}
G = (DD_*-a^2)v = \sum_i P_i\,\mathscr{G}_i(\alpha_i r),
\end{equation}
where the coefficients~$P_i$ in the expansion are left unspecified at this stage.
Having expressed~$G$ in this form, we next solve equation~\eqref{eq:7-299} as a
differential equation for~$v$ and arrange that it satisfies the remaining
boundary conditions, namely, that $v = 0$ for $r = 1$ and~$\eta$.  With~$v$
determined in this fashion, equation~\eqref{eq:7-291} will lead, as we shall presently
see, to a secular equation for~$T$.

\subsection*{(a) \textit{The characteristic equation}}

We shall now obtain the explicit form of the characteristic equation
for~$T$.  For this purpose it is convenient to write $\mathscr{G}_i(\alpha_i r)$ in the manner
\begin{equation}
\label{eq:7-300}
\mathscr{G}_i(\alpha_i r) = \mathfrak{u}_i(r) + v_i(r),
\end{equation}
where
\begin{align}
\label{eq:7-301}
\mathfrak{u}_i(r) &= A_i J_1(\alpha_i r) + B_i Y_1(\alpha_i r) \notag\\
\intertext{and}
v_i(r) &= C_i I_1(\alpha_i r) + D_i K_1(\alpha_i r).
\end{align}
Clearly,
\begin{equation}
\label{eq:7-302}
DD_*\,\mathfrak{u}_i = -\alpha_i^2\,\mathfrak{u}_i
\quad\text{and}\quad
DD_*\,v_i = +\alpha_i^2\,v_i.
\end{equation}
Making use of the relations~\eqref{eq:7-302}, we can readily verify that the general
solution of equation~\eqref{eq:7-299} is
\begin{equation}
\label{eq:7-303}
v = \sum_i P_i\!\left[
  p_i\,J_1(\alpha_i r) + q_i\,K_1(\alpha_i r)
  - \frac{\mathfrak{u}_i(r)}{\alpha_i^2+a^2}
  + \frac{v_i(r)}{\alpha_i^2-a^2}
\right],
\end{equation}
where the constants of integration $p_i$ and~$q_i$ are to be determined by
the boundary conditions on~$v$, namely, that~$v$ vanishes at $r = 1$ and~$\eta$.  These
latter conditions lead to the equations
\begin{equation}
\label{eq:7-304}
p_i\,J_1(a) + q_i\,K_1(a)
  = \frac{\mathfrak{u}_i(1)}{\alpha_i^2+a^2}
  - \frac{v_i(1)}{\alpha_i^2-a^2}
\end{equation}
and
\begin{equation}
p_i\,J_1(a\eta) + q_i\,K_1(a\eta)
  = \frac{\mathfrak{u}_i(\eta)}{\alpha_i^2+a^2}
  - \frac{v_i(\eta)}{\alpha_i^2-a^2}.
\end{equation}
In obtaining these equations from~\eqref{eq:7-303}, we have made use of the fact
that
\begin{equation}
\label{eq:7-305}
\mathfrak{u}_i(1) = -v_i(1) \quad\text{and}\quad
\mathfrak{u}_i(\eta) = -v_i(\eta).
\end{equation}
On solving equations~\eqref{eq:7-304}, we find:
\begin{align}
\label{eq:7-306}
p_i &= \frac{2\alpha_i^2}{\Delta(\alpha_i^2-a^2)}\bigl[
  +\mathfrak{u}_i(1)K_1(a\eta) - \mathfrak{u}_i(\eta)K_1(a)\bigr], \notag\\
q_i &= \frac{2\alpha_i^2}{\Delta(\alpha_i^2-a^2)}\bigl[
  -\mathfrak{u}_i(1)J_1(a\eta) + \mathfrak{u}_i(\eta)J_1(a)\bigr],
\end{align}
where
\begin{equation}
\label{eq:7-307}
\Delta = J_1(a\eta)K_1(a) - J_1(a)K_1(a\eta).
\end{equation}

Now substituting for~$v$ from equation~\eqref{eq:7-303} in equation~\eqref{eq:7-293}, we obtain
\begin{multline}
\label{eq:7-308}
\sum_i P_i\bigl[(\alpha_i^2+a^2)^2\mathfrak{u}_{ij} + (\alpha_i^2-a^2)^2v_{ij}\bigr] \\
= Ta^2\!\left(\tfrac{1}{2}-\kappa\right)\sum_j \sum_i P_i\!\left[
  p_i\,J_1(\alpha_i r) + q_i\,K_1(\alpha_i r)
  + \frac{\alpha_i^2(\mathfrak{u}_i + v_i)}{(\alpha_i^2+a^2)(\alpha_i^2-a^2)}
\right].
\end{multline}
Finally, multiplying this equation by $r(\mathfrak{u}_j + v_j)$ and integrating over the
range of~$r$, we obtain
\begin{multline}
\label{eq:7-309}
P_i(\alpha_i^2+a^2)N_i^2 + 2a^2\sum_j P_j \alpha_j^2 \Delta_i^{(j)} \\
= Ta^2\sum_j P_j\bigl[p_j \alpha \varepsilon I_j^{(+1)}(a)
  - I_j^{(-1)}(a)\bigr]
  + q_j\bigl[\alpha \varepsilon K_j^{(+1)}(a) - K_j^{(-1)}(a)\bigr] \\
  + \frac{\alpha_j^2}{(\alpha_j^2+a^2)}
    \bigl[\alpha N_{ij}^2 - \Delta_{ij}^{(-)}\bigr]
  - \frac{\alpha_j^2}{(\alpha_j^2-a^2)} M_{ij}\Biggr],
\end{multline}
where we have made use of the orthogonality relation
\begin{equation}
\label{eq:7-310}
\int (\mathfrak{u}_i+v_i)(\mathfrak{u}_j+v_j)\,dr = N_i^2\,\delta_{ij}.
\end{equation}
and introduced the abbreviations
\begin{align}
\label{eq:7-311}
\Delta_i^{(j)} &= \int (\mathfrak{u}_i-v_j)(\mathfrak{u}_i+v_j)a^{2+1}\,dr, \notag\\
I_j^{(\pm 1)}(a) &= \int J_1(ar)(\mathfrak{u}_j+v_j)a^{a\pm 1}\,dr, \notag\\
K_j^{(\pm 1)}(a) &= \int K_1(ar)(\mathfrak{u}_j+v_j)a^{a\pm 1}\,dr,
\end{align}
and
\begin{equation}
\label{eq:7-312}
M_{ij} = \int (\mathfrak{u}_i+v_i)(\mathfrak{u}_j+v_j)\frac{dr}{r}.
\end{equation}
Equation~\eqref{eq:7-309} leads to the secular equation
\begin{multline}
\label{eq:7-313}
\Biggl|N_i^2(\alpha_i^2+a^2)\,\delta_{ij}
  + 2a^2\alpha_j^2\Delta_i^{(j)} \\
-T\Bigl[p_j\,\alpha\,\varepsilon\,I_j^{(+1)}(a)
  - I_j^{(-1)}(a)\bigr]
  + q_j\bigl[\alpha\,\varepsilon\,K_j^{(+1)}(a) - K_j^{(-1)}(a)\bigr] \\
  - \frac{\alpha_j^2}{(\alpha_j^2+a^2)}
    \bigl[\alpha N_{ij}^2 - \Delta_{ij}^{(-)}\bigr]
  - \frac{\alpha_j^2}{(\alpha_j^2-a^2)}M_{ij}\Biggr|
= 0.
\end{multline}

Of the various matrix elements defined in equations~\eqref{eq:7-311} and~\eqref{eq:7-312},
those with the superscript~$(+1)$ in~\eqref{eq:7-311} can be evaluated explicitly.
The remaining matrix elements must be evaluated numerically.

\subsection*{(b) \textit{Numerical results for the case} $\eta = \frac{1}{2}$}

First, we may note that when $\eta = \frac{1}{2}$ the definitions of~$T$ and~$\kappa$ are
\begin{equation}
\label{eq:7-314}
\mathscr{T} = \frac{64}{9}\!\left(\frac{\Omega_1 R_1^2}{\nu}\right)^{\!2}
  (1-\mu)(1-4\mu)
\end{equation}
and
\begin{equation}
\label{eq:7-315}
\kappa = (1-4\mu)/(1-\mu).
\end{equation}
And Rayleigh's criterion for stability in this case is
\begin{equation}
\label{eq:7-316}
\mu > \tfrac{1}{4} \quad \text{(for stability).}
\end{equation}

\begin{table}[ht]
\centering
\caption{Taylor numbers for various assigned values of $\kappa$ and~$a$ ($\eta = \frac{1}{2}$)}
\label{tab:XXXIV}
\medskip
\small
\begin{tabular}{cccccc}
\hline
& & \textit{First} & \textit{Second} & \textit{Third} \\
$\kappa$ & $a$ & \textit{approxi-} & \textit{approxi-} & \textit{approxi-} \\
& & \textit{mation} & \textit{mation} & \textit{mation} \\
\hline
  & $6{\cdot}0$ & $1{\cdot}5530\times10^4$ & $1{\cdot}5470\times10^4$ & $1{\cdot}5370\times10^4$ \\
0 & $6{\cdot}2$ & $1{\cdot}5456\times10^4$ & $1{\cdot}5434\times10^4$ & $1{\cdot}5338\times10^4$ \\
  & $6{\cdot}4$ & $1{\cdot}5468\times10^4$ & $1{\cdot}5414\times10^4$ & $1{\cdot}5315\times10^4$ \\
\hline
    & $6{\cdot}0$ & $1{\cdot}5833\times10^4$ & $1{\cdot}5728\times10^4$ & \\
0.4 & $6{\cdot}2$ & $1{\cdot}5738\times10^4$ & $1{\cdot}5686\times10^4$ & $1{\cdot}5593\times10^4$ \\
    & $6{\cdot}4$ & $1{\cdot}5806\times10^4$ & $1{\cdot}5642\times10^4$ & $1{\cdot}5542\times10^4$ \\
\hline
    & $6{\cdot}0$ & $1{\cdot}6800\times10^4$ & & \\
0.6 & $6{\cdot}2$ & $1{\cdot}6604\times10^4$ & $1{\cdot}5642\times10^4$ & \\
    & $6{\cdot}4$ & $1{\cdot}6411\times10^4$ & $1{\cdot}6119\times10^4$ & \\
\hline
    & $5{\cdot}6$ & $1{\cdot}3890\times10^4$ & $1{\cdot}3386\times10^4$ & $1{\cdot}3110\times10^4$ \\
1.0 & $5{\cdot}8$ & $1{\cdot}3598\times10^4$ & $1{\cdot}3390\times10^4$ & $1{\cdot}3110\times10^4$ \\
    & $6{\cdot}0$ & $1{\cdot}3455\times10^4$ & $1{\cdot}3392\times10^4$ & \\
\hline
     & $5{\cdot}6$ & $1{\cdot}3468\times10^4$ & $1{\cdot}3492\times10^4$ & $1{\cdot}3352\times10^4$ \\
1.333 & $5{\cdot}8$ & $1{\cdot}5068\times10^4$ & $1{\cdot}3380\times10^4$ & \\
      & $6{\cdot}0$ & $1{\cdot}5316\times10^4$ & $1{\cdot}3307\times10^4$ & $1{\cdot}3045\times10^4$ \\
\hline
    & $5{\cdot}0$ & $1{\cdot}5035\times10^4$ & $1{\cdot}3807\times10^4$ & \\
1.6 & $5{\cdot}5$ & $1{\cdot}4961\times10^4$ & $1{\cdot}3007\times10^4$ & \\
    & $6{\cdot}0$ & $1{\cdot}1935\times10^4$ & $1{\cdot}1005\times10^4$ & $9{\cdot}8831\times10^3$ \\
    & $6{\cdot}5$ & $2{\cdot}0024\times10^4$ & $1{\cdot}4006\times10^4$ & \\
\hline
    & $7{\cdot}0$ & $7{\cdot}507\times10^4$ & & \\
1.8 & $7{\cdot}5$ & $7{\cdot}843\times10^4$ & $3{\cdot}0862\times10^4$ & \\
    & $8{\cdot}0$ & $8{\cdot}061\times10^4$ & & \\
\hline
    & $8{\cdot}4$ & & $2{\cdot}024$ & \\
1.9 & $8{\cdot}8$ & & $2{\cdot}065\times10^4$ & \\
    & $9{\cdot}2$ & & & \\
\hline
$\frac{9}{4}$ & $9{\cdot}6$ & & $5{\cdot}093\times10^4$ & $4{\cdot}305\times10^4$ \\
\hline
\end{tabular}
\end{table}

In Table~XXXIV the values of~$T$ obtained with the aid of the characteristic equation~\eqref{eq:7-313} in the different approximations (the `order' of the
approximation being the order of the determinant which is set equal
to zero in the determination of~$T$) are listed for various assigned values
of~$a$ and~$\kappa$.  For each value of~$\kappa$, values of~$a$ were chosen (by trial and
error) in the range in which~$T$ (as a function of~$a$) attains its minimum.
From an examination of this table, it appears that for $\kappa < 1{\cdot}6$ the
required values of~$T$ have been determined to within one per cent of the
true values.  For $\kappa = 1{\cdot}8$, $1{\cdot}9$, and~$2{\cdot}0$, the results obtained in the second
and the third approximations differ appreciably.  Nevertheless, it is
unlikely that even in these cases the results of the third approximation
are in error by more than 3~per cent: this estimate of the error does not
appear unreasonable when we observe that for $\kappa = 1{\cdot}4$, while the results
of the first and the second approximations differ by 20~per cent, the

\figurePlaceholder{72}{The critical Taylor number $T_c$ for the onset of instability as a function of $\kappa = (1-4\mu)/(1-\mu)$.  A scale of $\mu$ ($= \Omega_2/\Omega_1$) is also shown.}

results of the second and the third approximations do not differ by more
than 2~per cent.

In Table~XXXV the critical Taylor numbers (derived from the data
of Table~XXXIV) appropriate for the different values of~$\kappa$ are given;
the values of~$a$ at which these minimum Taylor numbers are attained are
also given.  The derived $(T_c,\,\kappa)$ and~$[a(T_c),\,\kappa]$-relationships are further
illustrated in Figs.~72 and~73.  From Fig.~72 it appears that in a $(\log T_c,\,\kappa)$-plot, the relationship is very nearly a linear one in the neighbourhood
of $\kappa = 2$.

Table~XXXV also includes the values
\begin{equation}
\label{eq:7-317}
\omega_R = \frac{\Omega_1 R_1^2}{\nu}
  = 0{\cdot}375\sqrt{\frac{T}{(1-\mu)(1-4\mu)}}
  = \frac{0{\cdot}375}{1-\mu}\sqrt{\frac{T}{\kappa}};
\end{equation}
and $\omega_R$ and $\omega_R$, are, therefore, the angular velocities $\Omega_1$ and~$\Omega_2$ measured in
the unit $\nu/R_1^2$.  In the $(\omega_R,\,\omega_R)$-plane, the locus determined by equations
\eqref{eq:7-317} separates the regions of stability from the regions of instability
(see Fig.~74).  In this plane, Rayleigh's criterion is represented by the
straight line
\begin{equation}
\label{eq:7-318}
\omega_R = \tfrac{1}{4}\omega_1 \quad \text{(Rayleigh's criterion).}
\end{equation}

\figurePlaceholder{73}{The variation of the wave number~$a$ (in the unit $1/R_2$) of the disturbance at which instability first sets in as a function of~$\kappa$ and~$\mu$.  The uncertainty in the determination of~$a$ is indicated by the height of the lines.}

\begin{table}[ht]
\centering
\caption{Critical Taylor numbers and related constants for various values of $\kappa$
  ($\eta = \frac{1}{2}$)}
\label{tab:XXXV}
\medskip
\small
\begin{tabular}{ccccccc}
\hline
$\mu$ & $\kappa$ & $a$ & $T_c$ & $\omega_R$ & $\omega_R$ \\
\hline
0       & $+0{\cdot}25$ & $6{\cdot}1$ & $1{\cdot}333\times10^4$ &         & \\
$10{\cdot}1$   & $+0{\cdot}12076$ & $6{\cdot}1$ & $1{\cdot}631\times10^4$ & $190{\cdot}3$ & $+42{\cdot}30$ \\
$10{\cdot}3$   & $-0{\cdot}19036$ & $6{\cdot}1$ & $1{\cdot}719\times10^4$ & $139{\cdot}3$ & $+49{\cdot}21$ \\
$10{\cdot}3$   & $+0{\cdot}31676$ & $6{\cdot}1$ & $1{\cdot}810\times10^4$ & $114{\cdot}0$ & \\
0.4     & $-0{\cdot}16667$ & $6{\cdot}2$ & $1{\cdot}954\times10^4$ & $69{\cdot}5$ & $+16{\cdot}38$ \\
$0{\cdot}6$ & $-0{\cdot}17647$ & $6{\cdot}2$ & $2{\cdot}864\times10^4$ & $81{\cdot}2$ & $-9{\cdot}73$ \\
\hline
$1{\cdot}0$ & $-0{\cdot}6$ & $5{\cdot}8$ & $1{\cdot}310\times10^4$ & $80{\cdot}4$ & \\
\hline
$1{\cdot}333$ & $-0{\cdot}425$ & $6{\cdot}0$ & $3{\cdot}335\times10^4$ & $66{\cdot}5$ & $-8{\cdot}310$ \\
$1{\cdot}5$ & & $6{\cdot}5$ & $6{\cdot}813\times10^4$ & $74{\cdot}15$ & \\
\hline
$1{\cdot}7$ & $-0{\cdot}30438$ & & $1{\cdot}377\times10^5$ & $81{\cdot}95$ & $+44{\cdot}60$ \\
$1{\cdot}8$ & & $7{\cdot}5$ & $1{\cdot}993\times10^5$ & $91{\cdot}10$ & $-23{\cdot}30$ \\
\hline
$1{\cdot}95$ & $-0{\cdot}30530$ & & $2{\cdot}435\times10^5$ & $97{\cdot}10$ & $-39{\cdot}42$ \\
$1{\cdot}9$ & $-0{\cdot}45573$ & & $2{\cdot}6$ & $1{\cdot}424\times10^5$ & $-44{\cdot}72$ \\
\hline
$1{\cdot}95$ & $-0{\cdot}49415$ & & $2{\cdot}550\times10^5$ & $109{\cdot}4$ & $-27{\cdot}71$ \\
$1{\cdot}99$ & $-0{\cdot}090$ & $9{\cdot}6$ & $4{\cdot}386\times10^5$ & & $-37{\cdot}79$ \\
\hline
\end{tabular}
\end{table}

Since (see the first entry in Table~XXXV)
\begin{equation}
\label{eq:7-319}
T_c \to 1{\cdot}553\times10^4 \quad \text{as } \mu \to \tfrac{1}{4}
  \text{ and } \kappa \to 0,
\end{equation}
it is clear that in the neighbourhood of the line~\eqref{eq:7-318}, as $\omega_R \to \infty$, the
asymptotic behaviour of the true $(\omega_R,\,\omega_R)$-locus is given by (cf.\ equation
\eqref{eq:7-317})
\begin{equation}
\label{eq:7-320}
\omega_2 = 0{\cdot}375\sqrt{\frac{1{\cdot}553\times10^4}{\kappa(1-4\mu)}}
  = \frac{53{\cdot}81}{\sqrt{\kappa}} \quad (\mu \to \tfrac{1}{4}).
\end{equation}

\figurePlaceholder{74}{The regions of stability and instability in the $(\omega_R,\,\omega_R)$-plane; $\omega_1$ and~$\omega_2$ are the angular velocities of the inner and the outer cylinder in the units $\nu/R_1^2$ and~$\nu/R_2^2$, respectively.}

Finally, in Figs.~75 and~76, the profiles of the velocity in the radial
direction and the corresponding cell patterns at marginal stability are
shown for two typical cases ($\mu = 0$ and $\mu = -4/11$).

\figurePlaceholder{75}{The velocity profile~(a) and the cell pattern~(b) at the onset of instability for the case $\kappa = 1$, $\mu = 0$, and $a = 6{\cdot}4$.  The stream function $\psi$ (or $v_z \cos ax$) has been normalized to unity and the cell pattern is drawn symmetrically about $x = 0$.  (The unit of length is the radius of the outer cylinder.)}

\figurePlaceholder{76}{The velocity profile~(a) and the cell pattern~(b) at the onset of instability for the case $\kappa = 1{\cdot}8$, $\mu = -4/11$, and $a = 7{\cdot}6$.  The stream function $\psi$ (or $v_z \cos ax$) has been normalized to unity and the cell pattern is drawn symmetrically about $x = 0$.  (The unit of length is the radius of the outer cylinder.)}

%% §74 %%
\section{Experiments on the stability of viscous flow between rotating cylinders}\label{sec:7-74}

% PLACEHOLDER: Section 74 content

%% BIBLIOGRAPHICAL NOTES %%
\bigskip
\begin{center}
\textsc{BIBLIOGRAPHICAL NOTES}
\end{center}

% PLACEHOLDER: Bibliographical Notes content
